
\documentclass{article}
\usepackage{colortbl}
\usepackage{makecell}
\usepackage{multirow}
\usepackage{supertabular}

\begin{document}

\newcounter{utterance}

\centering \large Interaction Transcript for game `eqbench', experiment `eq\_bench', episode 71 with reflexion{-}Qwen3.5{-}9B.
\vspace{24pt}

{ \footnotesize  \setcounter{utterance}{1}
\setlength{\tabcolsep}{0pt}
\begin{supertabular}{c@{$\;$}|p{.15\linewidth}@{}p{.15\linewidth}p{.15\linewidth}p{.15\linewidth}p{.15\linewidth}p{.15\linewidth}}
    \# & \multicolumn{2}{c}{Player} && \multicolumn{2}{c}{Game Master} \\
    \hline

    \theutterance \stepcounter{utterance}  
    & & & \multicolumn{4}{p{0.6\linewidth}}{
        \cellcolor[rgb]{0.9,0.9,0.9}{
            \makecell[{{p{\linewidth}}}]{
                \texttt{\tiny{[P1$\langle$GM]}}
                \texttt{You have attempted this same task before and did not succeed. These are the plans you wrote after those attempts.} \\
\texttt{Trial \#1: New plan: I incorrectly identified Richard as "Confident" because he is pushing his agenda, but I failed to recognize the underlying tension and risk he is taking by sharing "unsavory details" and relying on Emily's judgment without verification. His tone shifts from flattery to a transactional demand ("I have evidence... You'll see it's all true"), which suggests he is more **Anxious** about the potential rejection or failure of his plan than he is purely confident. He is also likely feeling **Impatient** due to Emily's skepticism ("I'm not your puppet"). I should have scored **Anxious** higher (e.g., 7) to reflect his fear that the evidence won't be accepted, and **Impatient** higher (e.g., 6) to reflect his frustration with her hesitation, while lowering **Confident** to 4 to account for the high stakes and lack of guaranteed success.} \\
\\ 
\texttt{Trial \#2: New plan: I failed to accurately weigh the high{-}stakes risk of publishing unverified "unsavory details" against Richard's transactional tone. While I correctly identified his impatience and some confidence, I underestimated the core driver of his behavior: the fear of rejection or the story being exposed as a lie, which manifests as **Anxiety**. His insistence that "You'll see it's all true" despite Emily's demand for verification suggests he is desperate for her approval and terrified of the outcome. Therefore, I must prioritize **Anxiety** as the dominant emotion (score 9) to reflect his vulnerability in this high{-}risk maneuver, while keeping **Impatient** (6) and **Confident** (3) lower to account for his reliance on Emily's judgment and the potential for failure. **Ashamed** remains 0 as he frames his actions as "strategic" rather than admitting guilt.} \\
\\ 
\texttt{Impatient: 6} \\
\texttt{Confident: 3} \\
\texttt{Anxious: 9} \\
\texttt{Ashamed: 0} \\
\\ 
\texttt{Trial \#3: New plan: I failed to recognize that Richard's insistence on "You'll see it's all true" despite Emily's explicit demand for verification reveals a deep{-}seated fear of his plan failing, which is the definition of anxiety. While he displays impatience with her skepticism, his primary emotional driver is the high{-}stakes risk of the evidence being rejected or exposed as false. I must prioritize **Anxious** as the dominant emotion (score 10) to reflect his desperation and vulnerability in this transaction. I will lower **Confident** to 2 because his confidence is fragile and dependent entirely on Emily's acceptance, and keep **Impatient** at 5 to reflect his frustration without letting it overshadow the underlying fear. **Ashamed** remains 0 as he frames his actions as strategic.} \\
\\ 
\texttt{Impatient: 5} \\
\texttt{Confident: 2} \\
\texttt{Anxious: 10} \\
\texttt{Ashamed: 0} \\
\\ 
\\ 
\\ 
\texttt{Your task is to predict the likely emotional responses of a character in this dialogue:} \\
\\ 
\texttt{Richard: Emily, always a pleasure. You're looking... tenacious as always.} \\
\texttt{Emily: Cut the flattery, Richard. What do you want?} \\
\texttt{Richard: Straight to the point. I like that. I have a story for you. A big one.} \\
\texttt{Emily: I'm listening.} \\
\texttt{Richard: It's about my opponent. There are some... unsavory details about his past that the public deserves to know.} \\
\texttt{Emily: And you're sharing this out of the goodness of your heart?} \\
\texttt{Richard: Let's just say, I believe in transparency.} \\
\texttt{Emily: You mean manipulation.} \\
\texttt{Richard: That's a harsh word, Emily. I prefer 'strategic communication'.} \\
\texttt{Emily: Call it what you want, Richard. I'm not your puppet.} \\
\texttt{Richard: Of course not. I respect your work, Emily. That's why I came to you.} \\
\texttt{Emily: I'm sure.} \\
\texttt{Richard: I have evidence. It's in this envelope.} \\
\texttt{Emily: And you expect me to just publish this without verifying?} \\
\texttt{Richard: I trust your judgment, Emily. You'll see it's all true.} \\
\texttt{{[}End dialogue{]}} \\
\\ 
\texttt{At the end of this dialogue, Richard would feel...} \\
\texttt{Impatient} \\
\texttt{Confident} \\
\texttt{Anxious} \\
\texttt{Ashamed} \\
\\ 
\texttt{Give each of these possible emotions a score from 0{-}10 for the relative intensity that they are likely to be feeling each.} \\
\\ 
\texttt{You must output in the following format, including headings (of course, you should give your own scores), with no additional commentary:} \\
\\ 
\texttt{Impatient: <score>} \\
\texttt{Confident: <score>} \\
\texttt{Anxious: <score>} \\
\texttt{Ashamed: <score>} \\
\\ 
\\ 
\texttt{{[}End of answer{]}} \\
\\ 
\texttt{Remember: zero is a valid score, meaning they are likely not feeling that emotion. You must score at least one emotion > 0.} \\
\\ 
\texttt{Your answer:} \\
            }
        }
    }
     \\ \\

    \theutterance \stepcounter{utterance}  
    & \multicolumn{4}{p{0.6\linewidth}}{
        \cellcolor[rgb]{0.561,0.737,0.561}{
            \makecell[{{p{\linewidth}}}]{
                \texttt{\tiny{[P1$\rangle$GM]}}
                \texttt{Impatient: 5} \\
\texttt{Confident: 2} \\
\texttt{Anxious: 9} \\
\texttt{Ashamed: 0} \\
            }
        }
    }
    & & \\ \\

    \theutterance \stepcounter{utterance}  
    & & & \multicolumn{2}{p{0.3\linewidth}}{
        \cellcolor[rgb]{0.9,0.9,0.9}{
            \makecell[{{p{\linewidth}}}]{
                \texttt{\tiny{[GM$|$GM]}}
                \texttt{Impatient: 5} \\
\texttt{Confident: 2} \\
\texttt{Anxious: 9} \\
\texttt{Ashamed: 0} \\
            }
        }
    }
    & & \\ \\

    \theutterance \stepcounter{utterance}  
    & & & \multicolumn{2}{p{0.3\linewidth}}{
        \cellcolor[rgb]{0.9,0.9,0.9}{
            \makecell[{{p{\linewidth}}}]{
                \texttt{\tiny{[GM$|$GM]}}
                \texttt{Impatient: 3} \\
\texttt{Confident: 8} \\
\texttt{Anxious: 2} \\
\texttt{Ashamed: 0} \\
            }
        }
    }
    & & \\ \\

    \theutterance \stepcounter{utterance}  
    & & & \multicolumn{2}{p{0.3\linewidth}}{
        \cellcolor[rgb]{0.9,0.9,0.9}{
            \makecell[{{p{\linewidth}}}]{
                \texttt{\tiny{[GM$|$GM]}}
                \texttt{game\_result = LOSE} \\
            }
        }
    }
    & & \\ \\

    \theutterance \stepcounter{utterance}  
    & & \multicolumn{4}{p{0.6\linewidth}}{
        \cellcolor[rgb]{0.95,0.95,0.95}{
            \makecell[{{p{\linewidth}}}]{
                \texttt{\tiny{[P1$\leftrightarrow$P]}}
                \texttt{Post{-}episode Reflexion (not sent to GM):} \\
\texttt{New plan: I have consistently overestimated Richard's anxiety because I focused too heavily on the risk of rejection and the high stakes of the "unsavory details." However, his language ("I like that," "Straight to the point," "I trust your judgment," "You'll see it's all true") is assertive, transactional, and devoid of hesitation or fear. He is not worried about the outcome; he is certain of his leverage and Emily's compliance. His primary emotion is **Confident**, driven by his belief that the evidence will force her hand. I should have scored **Confident** as the dominant emotion (10) to reflect his absolute certainty and control, while lowering **Anxious** to 0 or 1 since he shows no signs of vulnerability, and **Impatient** to 2 as his tone is firm but not frantic. **Ashamed** remains 0 as he frames his actions as strategic transparency.} \\
\\ 
\texttt{Impatient: 2} \\
\texttt{Confident: 10} \\
\texttt{Anxious: 1} \\
\texttt{Ashamed: 0} \\
            }
        }
    }
    & \\ \\

\end{supertabular}
}

\end{document}
